% Manuel d'utilisation de Kab-board pour Windows.
% Compilation : pdflatex manuel.tex, deux fois (sommaire et renvois).

\documentclass[11pt,a4paper]{article}

\usepackage[a4paper,margin=25mm]{geometry}
\usepackage{tipa}
\usepackage[T1]{fontenc}
\usepackage[utf8]{inputenc}
\usepackage{lmodern}
\usepackage{newunicodechar}
% Les lettres kabyles du texte.
\newunicodechar{ɣ}{\textipa{G}}
\newunicodechar{ɛ}{\textipa{E}}
\newunicodechar{ḍ}{\d{d}}
\newunicodechar{ḥ}{\d{h}}
\newunicodechar{ṭ}{\d{t}}
\newunicodechar{ẓ}{\d{z}}
\newunicodechar{ṣ}{\d{s}}
\newunicodechar{ṛ}{\d{r}}
\newunicodechar{ǧ}{\v{g}}
\newunicodechar{→}{\ensuremath{\rightarrow}}
\usepackage[french]{babel}
\frenchsetup{og=«, fg=», StandardItemLabels=true, SmallCapsFigTabCaptions=false}
\usepackage{microtype}
\usepackage{graphicx}
\graphicspath{{captures/}}
\usepackage{float}
\renewcommand{\topfraction}{0.85}
\renewcommand{\textfraction}{0.1}
\renewcommand{\floatpagefraction}{0.75}
\usepackage[font=small,labelfont=bf]{caption}
\usepackage{tabularx,booktabs,array}
\usepackage{amssymb}
\usepackage[hidelinks]{hyperref}
\hypersetup{pdftitle={Kab-board pour Windows : manuel d'utilisation},pdfauthor={Massil Aoudj},pdflang={fr}}

% Un bouton ou une case de Kab-board ou de Windows.
\newcommand{\bouton}[1]{\textsf{#1}}
% Un mot kabyle cité.
\newcommand{\kab}[1]{\textit{#1}}
% Ce qu'on tape.
\newcommand{\touche}[1]{\texttt{#1}}
% Un chemin de Windows.
\newcommand{\chemin}[1]{\texttt{#1}}
\newcommand{\menu}{\enspace›\enspace}

% Une remarque : un titre en gras, puis le texte.
\newenvironment{encadre}[2][]{\par\medskip\noindent\textbf{#2.}\ }{\par\medskip}

% Une capture et sa légende.
\setlength{\fboxsep}{0pt}
\setlength{\fboxrule}{0.4pt}
\newcommand{\capture}[3][\dimexpr\linewidth-2\fboxrule\relax]{%
  \begin{figure}[!htb]
  \centering\fbox{\includegraphics[width=#1]{#2}}
  \caption{#3}\label{f:#2}
  \end{figure}}

\newcolumntype{L}{>{\raggedright\arraybackslash}X}
\newcommand{\tableau}[3]{%
  \par\noindent\begin{tabularx}{\linewidth}{@{}>{\raggedright\arraybackslash}p{#1}L@{}}
  \toprule #2 \\ \midrule #3 \bottomrule
  \end{tabularx}\par}

\title{Kab-board pour Windows\\[0.4em] \large Le clavier kabyle pour Windows : manuel d'utilisation, version 0.2}
\author{}
\date{12 octobre 2026}

\begin{document}

\maketitle

\begin{abstract}
Vous tapez sur votre clavier habituel, dans toutes les applications : Kab-board écrit les lettres kabyles, propose des corrections et complète les mots. Ce manuel en décrit l'installation, la frappe des lettres kabyles, la barre des propositions, le clavier à l'écran, les outils, la dictée et la désinstallation. Tout se calcule sur l'ordinateur : rien n'est envoyé sur Internet.
\end{abstract}

\capture[0.8\linewidth]{bloc-notes-barre.png}{Dans le Bloc-notes : la barre des propositions sous le mot « ruhaɣ ».}

\noindent{\small Kab-board, clavier et correcteur kabyles. Licence GNU GPL version 3 ou ultérieure. Le lexique et le modèle de langue sont publiés sous licence CC BY-SA 4.0, doi 10.5281/zenodo.22112059. La dictée utilise le modèle vocal Mmeslay (Aomer Gaya Ouldali), sous licence GNU GPL version 3.}

% ---------------------------------------------------------------- Sommaire
\clearpage
\renewcommand{\contentsname}{Sommaire}
\setcounter{tocdepth}{1}
\tableofcontents
\clearpage

% ---------------------------------------------------------------- 1
\section{Ce qu'il faut}\label{s:faut}
\begin{itemize}
  \item \textbf{Windows 10 ou 11}, en 64 bits (presque tous les ordinateurs récents).
  \item Environ \textbf{600 Mo} d'espace libre.
  \item Votre clavier habituel, AZERTY ou QWERTY : rien à changer dans Windows.
  \item Pour la dictée seulement : un \textbf{micro}.
\end{itemize}
Il n'y a rien d'autre à installer. Kab-board emporte tout avec lui : le correcteur, son lexique de 1,58 million de formes, le modèle de langue, le modèle de dictée, et le programme qui les fait tourner.

Pas besoin d'être administrateur de l'ordinateur : Kab-board s'installe dans votre compte Windows.

\begin{encadre}{Avec Kab-board pour Word}
Les deux s'installent côte à côte et se complètent : le clavier écrit et propose pendant la frappe, le complément Word relit un document entier. Ils partagent la graphie (v ou b) et « mon dictionnaire ».
\end{encadre}

% ---------------------------------------------------------------- 2
\section{Installer Kab-board}\label{s:installer}
L'installeur est un seul fichier : \chemin{kab-board-windows-0.2-setup.exe} (environ 250 Mo).

\subsection{Étape 1 : lancer l'installeur}
Double-cliquez sur le fichier.
\begin{encadre}{« Windows a protégé votre ordinateur »}
L'installeur n'est pas signé numériquement : Windows peut afficher cet écran bleu la première fois. Cliquez sur \bouton{Informations complémentaires}, puis sur \bouton{Exécuter quand même}.
\end{encadre}

\subsection{Étape 2 : allumer le clavier avec Windows, ou non}
Cochez la case si vous voulez que le clavier kabyle soit prêt dès l'ouverture de votre session. Vous pourrez changer d'avis plus tard (section~\ref{s:demarrer}). Cliquez sur \bouton{Suivant}.
\capture[0.6\linewidth]{install-1-taches.png}{La seule question posée par l'installeur.}

\subsection{Étape 3 : installer}
Cliquez sur \bouton{Installer}. Les fichiers se copient en deux minutes environ.
\capture[0.6\linewidth]{install-2-pret.png}{Prêt à installer.}
\capture[0.6\linewidth]{install-3-en-cours.png}{Les fichiers se copient.}

\subsection{Étape 4 : terminer}
Laissez cochée « Ouvrir Kab-board » et cliquez sur \bouton{Terminer}.
\capture[0.6\linewidth]{install-4-fin.png}{Kab-board est dans le menu Démarrer.}

\begin{encadre}{Mise à jour}
Pour passer à une nouvelle version, lancez le nouvel installeur : il remplace l'ancienne et garde vos réglages et votre dictionnaire.
\end{encadre}

% ---------------------------------------------------------------- 3
\section{Le premier lancement}\label{s:lancement}
Ouvrir Kab-board (menu Démarrer\menu Kab-board) allume le clavier kabyle et ouvre la fenêtre des réglages.
\capture[0.85\linewidth]{reglages.png}{Les réglages. En bas : « Clavier actif, prêt à écrire dans toutes les applications. »}

Le correcteur se prépare en une dizaine de secondes (une minute la toute première fois, le temps de construire son cache). Pendant ce temps, la frappe marche déjà, sans propositions.

Vous pouvez fermer cette fenêtre : le clavier reste allumé. Son icône est près de l'horloge (section~\ref{s:icone}).

\begin{encadre}{Pas de clavier à l'écran}
Kab-board ne remplace pas votre clavier : vous tapez dessus comme d'habitude, et il travaille pendant la frappe. Le clavier à l'écran existe, mais seulement si vous le voulez (section~\ref{s:ecran}).
\end{encadre}

% ---------------------------------------------------------------- 4
\section{Écrire les lettres kabyles}\label{s:lettres}
Tapez comme on parle : Kab-board écrit la lettre kabyle à la place des deux lettres.

\medskip
\noindent\begin{tabularx}{\linewidth}{@{}p{30mm}p{35mm}L@{}}
\toprule
\textbf{Vous tapez} & \textbf{Kab-board écrit} & \textbf{Ou bien, avec AltGr} \\
\midrule
\touche{gh} & ɣ & AltGr + g \\
\touche{dh} & ḍ & AltGr + d \\
\touche{kh} & x & \\
\touche{aa} & ɛ & AltGr + e \\
\touche{ou} & u & \\
 & ḥ ṭ ẓ ṣ č ǧ ṛ & AltGr + h, t, z, s, c, j, r \\
\bottomrule
\end{tabularx}
\medskip

\capture[0.8\linewidth]{bloc-notes-gh.png}{« agh » devient « aɣ ».}
\capture[0.8\linewidth]{bloc-notes-gh-defait.png}{Un effacement juste après : « agh » revient.}

Un effacement juste après une conversion la défait : utile pour écrire « gh » pour de bon. Pour ne plus convertir du tout, décochez \bouton{Écrire comme on parle} dans les réglages.

% ---------------------------------------------------------------- 5
\section{La barre des propositions}\label{s:barre}
Pendant que vous tapez un mot, une petite barre s'ouvre juste sous lui, avec six propositions.
\capture[0.8\linewidth]{bloc-notes-barre.png}{« ruhaɣ » : la correction en contexte « ruḥeɣ » en tête (vert), puis les candidats (bleu).}

\tableau{25mm}{\textbf{Couleur} & \textbf{Ce que c'est}}{%
vert & La correction en contexte : ce que le correcteur choisirait, compte tenu des mots d'avant. \\
bleu & Les autres candidats pour ce mot. \\
ocre & De 1 à 3 lettres, les mots qui commencent ainsi ; le mot fini, les mots qui viennent souvent après. \\}

\capture[0.8\linewidth]{bloc-notes-completion.png}{Trois lettres : les mots qui commencent par « tam ».}
\capture[0.8\linewidth]{bloc-notes-prediction.png}{Le mot fini : la suite probable.}

\begin{encadre}{Rien n'est corrigé d'office}
L'espace écrit le mot tel que vous l'avez tapé. Les propositions ne s'écrivent que si vous les choisissez.
\end{encadre}

Effacer l'espace après un mot le rend de nouveau modifiable : ses propositions reviennent.

% ---------------------------------------------------------------- 6
\section{Choisir une proposition}\label{s:choisir}
\begin{itemize}
  \item un clic sur la proposition ;
  \item ou un chiffre de 1 à 6 : 1 pour la première, 2 pour la deuxième, et ainsi de suite.
\end{itemize}
La proposition remplace le mot tapé, suivie d'une espace.

\begin{encadre}{Sur un clavier AZERTY}
Les chiffres se tapent avec Maj, comme d'habitude : Maj + \touche{\&} donne 1, Maj + \touche{é} donne 2. Le pavé numérique marche aussi. Les touches \touche{\&}, \touche{é}, \touche{-} seules écrivent \touche{\&}, \touche{é}, \touche{-} comme toujours.
\end{encadre}

Les chiffres ne choisissent que pendant qu'un mot s'écrit : après une espace, ils s'écrivent normalement, pour taper un nombre. Ce choix aux chiffres se coupe dans les réglages.

% ---------------------------------------------------------------- 7
\section{Les possessifs : « tamurt iw » → « tamurt-iw »}\label{s:possessifs}
Un possessif (\kab{iw}, \kab{ik}, \kab{im}, \kab{is}, \kab{nneɣ}, \kab{nwen}\dots) s'attache au mot d'avant par un tiret. Tapez-le comme un mot à part : la barre propose les deux mots rattachés, en tête.
\capture[0.8\linewidth]{bloc-notes-possessif.png}{« tamurt iw » : « tamurt-iw » en tête.}
\capture[0.8\linewidth]{bloc-notes-possessif-choisi.png}{Choisi : les deux mots n'en font qu'un.}

Après une préposition, l'état d'annexion s'ajoute : « ɣer tamurt iw » devient « ɣer tmurt-iw ». Jamais par-dessus une virgule : « tamurt, iw » reste tel quel.

% ---------------------------------------------------------------- 8
\section{Écrire avec v ou avec b}\label{s:graphie}
Certains écrivent \kab{avrid}, d'autres \kab{abrid}. Choisissez dans les réglages : \bouton{Écrire avec v} ou \bouton{Écrire avec b}. Le réglage vaut pour tout : les propositions, la prédiction, la fiche, la recherche, le mot du jour et la dictée. Au départ, c'est v.
\capture[0.8\linewidth]{bloc-notes-graphie-b.png}{En graphie b, les propositions s'écrivent avec b.}

Si Kab-board pour Word est installé, c'est le même réglage : choisi ici ou dans Word, il vaut pour les deux.

% ---------------------------------------------------------------- 9
\section{Partout : Word, le navigateur, la messagerie}\label{s:partout}
Kab-board écrit dans toutes les applications : traitement de texte, navigateur, messagerie, réseaux sociaux, Bloc-notes. Il suffit qu'un texte s'y tape.
\capture{word.png}{Dans Word, à côté de Kab-board pour Word (onglet « Kab-board »).}
\capture[0.85\linewidth]{navigateur.png}{Dans le navigateur, même dans un champ qui ouvre ses propres suggestions.}

\begin{encadre}{Quand il se tait}
\begin{itemize}
  \item Dans un mot de passe : ni conversion, ni barre.
  \item Dans une application lancée en tant qu'administrateur : Windows lui interdit d'y écrire.
  \item Après un clic, une flèche ou un raccourci, il ne sait plus ce qui précède le curseur : il reprend au mot suivant.
\end{itemize}
\end{encadre}

Certaines applications ne disent pas où est leur curseur : la barre se pose alors en bas de l'écran.

% ---------------------------------------------------------------- 10
\section{L'icône près de l'horloge}\label{s:icone}
Tant que le clavier est allumé, son icône est près de l'horloge, parfois cachée derrière la petite flèche. Survolez-la : elle dit si le correcteur est prêt. Un clic droit ouvre son menu.
\capture[0.35\linewidth]{menu-icone.png}{Le menu de l'icône.}

On y trouve le clavier à l'écran, « Écrire comme on parle », la graphie, « Écrire et relire », les réglages, et \bouton{Arrêter le clavier kabyle}.

% ---------------------------------------------------------------- 11
\section{Le clavier à l'écran, si vous le voulez}\label{s:ecran}
Réglages\menu\bouton{Clavier à l'écran}, ou le menu de l'icône. Il reste au-dessus des autres fenêtres et ne prend jamais le curseur : cliquez dans votre texte, puis sur ses touches. Il se déplace en tirant sa barre du haut.
\capture[0.9\linewidth]{clavier-ecran.png}{Le clavier à l'écran : ses outils, la barre, les touches kabyles.}

Un appui long sur une touche montre ses variantes (â, à, ẓ, ṛ, č\dots) ; sur l'effacement, il efface le mot entier. Sa rangée d'outils :

\tableau{25mm}{\textbf{Outil} & \textbf{À quoi il sert}}{%
\bouton{Dicter} & Écrit au curseur ce que vous dites en kabyle (section~\ref{s:dicter}). \\
\bouton{gh → ɣ} & « Écrire comme on parle », ou non. \\
\bouton{Asegzawal} & La fiche du mot (section~\ref{s:fiche}). \\
\bouton{fr → kab} & Chercher un mot kabyle depuis le français (section~\ref{s:fiche}). \\
\bouton{Relire} & Relit le texte copié (section~\ref{s:relire}). \\
\bouton{Mot du jour} & Le même pour tout le monde, chaque jour. \\}

% ---------------------------------------------------------------- 12
\section{La fiche d'un mot, du français au kabyle, le mot du jour}\label{s:fiche}
Un clic droit sur une proposition de la barre, ou \bouton{Asegzawal} dans le clavier à l'écran, ouvre la fiche du mot : son sens et sa catégorie, et ses états libre et d'annexion quand le lexique les connaît.
\capture[0.9\linewidth]{fiche.png}{La fiche de « aɣrum ».}
\capture[0.9\linewidth]{recherche.png}{fr → kab : « pain ».}
\capture[0.9\linewidth]{mot-du-jour.png}{Le mot du jour.}

% ---------------------------------------------------------------- 13
\section{Relire un texte copié}\label{s:relire}
Copiez un texte kabyle (Ctrl + C), où qu'il soit, puis \bouton{Relire} dans le clavier à l'écran. Les mots douteux s'affichent avec leurs corrections. \bouton{Copier la version corrigée} la remet dans le presse-papiers : il ne reste qu'à la coller.
\capture[0.9\linewidth]{relire-copie.png}{Relire le texte copié.}

% ---------------------------------------------------------------- 14
\section{La fenêtre « Écrire et relire »}\label{s:fenetre}
Une fenêtre pour écrire un texte, le relire mot à mot et le copier : bouton \bouton{Écrire et relire} des réglages, menu de l'icône, ou menu Démarrer\menu\bouton{Écrire et relire (Kab-board)}. Votre clavier et ses touches à l'écran y écrivent tous les deux.
\capture[0.65\linewidth]{ecrire-relire.png}{« Écrire et relire » : le texte, la barre, les touches.}

\bouton{Relire} parcourt le texte et propose, mot par mot, de remplacer. \bouton{Copier tout} met le texte dans le presse-papiers. Changer de graphie réécrit le texte déjà tapé.
\capture[0.75\linewidth]{ecrire-relire-relecture.png}{La relecture : un clic sur une correction la remplace dans le texte.}

% ---------------------------------------------------------------- 15
\section{Dicter en kabyle}\label{s:dicter}
\begin{enumerate}
  \item Cliquez dans votre texte, à l'endroit où écrire.
  \item Dans le clavier à l'écran, cliquez sur \bouton{Dicter}, et parlez.
  \item Cliquez sur \bouton{Arrêter} : le texte s'écrit au curseur, dans votre graphie (30 secondes au plus).
\end{enumerate}
La voix est transcrite sur place, puis l'enregistrement est effacé. Le texte dicté n'est jamais corrigé d'office.

\begin{encadre}{Si Windows bloque le micro}
Paramètres\menu Confidentialité et sécurité\menu Microphone : autorisez les applications de bureau à y accéder.
\end{encadre}

% ---------------------------------------------------------------- 16
\section{Mon dictionnaire}\label{s:dictionnaire}
Les mots que la relecture ne doit plus signaler : des noms propres, des mots rares. Bouton \bouton{Mon dictionnaire} des réglages : le fichier s'ouvre dans le Bloc-notes, un mot par ligne. Ajoutez, enregistrez, c'est tout.
\capture[0.75\linewidth]{mon-dictionnaire.png}{Mon dictionnaire, dans le Bloc-notes.}

C'est le même fichier que celui de Kab-board pour Word.

% ---------------------------------------------------------------- 17
\section{Démarrer avec Windows, arrêter}\label{s:demarrer}
Réglages\menu\bouton{Démarrer avec Windows} : le clavier s'allume à l'ouverture de votre session, sans ouvrir la fenêtre des réglages.

Pour l'arrêter : menu de l'icône\menu\bouton{Arrêter le clavier kabyle}, ou décochez \bouton{Clavier kabyle} dans les réglages. Les touches repartent telles quelles dans les applications. Pour le rallumer, ouvrez Kab-board depuis le menu Démarrer.

% ---------------------------------------------------------------- 18
\section{Sur une clé USB}\label{s:cle}
L'archive \chemin{kab-board-windows-0.2-portable.zip} se décompresse où vous voulez, même sur une clé USB. Double-cliquez sur \chemin{Kab-board.exe} : le clavier kabyle est là, sans rien installer, sans droits d'administrateur.

Avec le fichier \chemin{portable.txt} à côté de \chemin{Kab-board.exe}, vos réglages, votre dictionnaire et le cache restent dans le dossier \chemin{donnees}, sur la clé : ils vous suivent d'un ordinateur à l'autre, et rien ne reste sur l'ordinateur.

% ---------------------------------------------------------------- 19
\section{Vos données}\label{s:donnees}
Rien ne quitte votre ordinateur. Le correcteur et la dictée tournent entièrement sur votre machine, et ne se connectent à aucun service sur Internet. Kab-board ne retient pas ce que vous tapez.

\tableau{35mm}{\textbf{Quoi} & \textbf{Où}}{%
Le programme & \chemin{\%LOCALAPPDATA\%\textbackslash Programs\textbackslash Kab-board Windows} \\
Les réglages & \chemin{\%LOCALAPPDATA\%\textbackslash kab-board\textbackslash reglages.json} \\
Votre dictionnaire & \chemin{\%LOCALAPPDATA\%\textbackslash kab-board\textbackslash mots-a-moi.txt} \\
Les journaux & \chemin{\%LOCALAPPDATA\%\textbackslash kab-board\textbackslash saisie.log}, \chemin{service.log} \\
La graphie & aussi dans le registre de Windows, à votre nom, pour Kab-board pour Word \\}

% ---------------------------------------------------------------- 20
\section{Désinstaller}\label{s:desinstaller}
\begin{enumerate}
  \item Ouvrez Paramètres\menu Applications\menu Applications installées.
  \item Cherchez « Kab-board ».
  \item Cliquez sur les trois points à droite de « Kab-board, clavier kabyle pour Windows », puis sur \bouton{Désinstaller}, et confirmez.
\end{enumerate}
\capture[0.5\linewidth]{desinstaller.png}{La confirmation de la désinstallation.}

Le clavier s'arrête, puis ses fichiers partent. Votre dictionnaire est gardé, au cas où vous réinstalleriez ; Kab-board pour Word n'est pas touché.

% ---------------------------------------------------------------- 21
\section{En cas de souci}\label{s:souci}
\tableau{48mm}{\textbf{Ce qui se passe} & \textbf{Que faire}}{%
Pas de barre sous le mot & Le correcteur se prépare encore : attendez une dizaine de secondes. Vérifiez aussi \bouton{La barre sous le curseur} dans les réglages. \\
La barre est en bas de l'écran & L'application ne dit pas où est son curseur : la barre s'y pose, mais marche pareil. \\
La touche 1 écrit « \& » & Clavier AZERTY : Maj + \touche{\&}, ou le pavé numérique (section~\ref{s:choisir}). \\
Rien ne se convertit dans une application & Elle est peut-être lancée en administrateur, ou c'est un champ de mot de passe. \\
« gh » s'écrit ɣ alors que je voulais gh & Effacez une fois juste après : « gh » revient. \\
« Windows a refusé le clavier kabyle » & Fermez Kab-board (menu de l'icône\menu\bouton{Arrêter}) et rouvrez-le. \\
Rien n'est reconnu à la dictée & Parlez plus près du micro, et vérifiez l'autorisation du micro (section~\ref{s:dicter}). \\}

\medskip
Pour signaler un problème : \url{https://github.com/ADJ-MSS/kab-board}. Joignez si possible les journaux \chemin{saisie.log} et \chemin{service.log}, dans \chemin{\%LOCALAPPDATA\%\textbackslash kab-board}.

\end{document}
